\chapter{Poisson Processes}

\begin{center}\itshape
Buses arrive on average every ten minutes, at random. You reach the stop at a random time.\\
Your expected wait is not five minutes but ten, and the gap you fall into lasts twenty on average.
\end{center}

In 1898 Ladislaus Bortkiewicz published a small book called \emph{The Law of Small Numbers}. Its most famous table counts, for twenty years and
fourteen corps of the Prussian cavalry, the soldiers killed each year by the kick of a horse. Most corps lost no one in most years, a few lost one or
two, and one lost four; and the counts followed, with uncanny precision, the law that Sim\'eon Denis Poisson had written down sixty years earlier for
rare events. This chapter explains why rare, independent events always obey Poisson, and turns his law into a process in time.

Calls reaching a switchboard, requests hitting a web server, decays of radioactive atoms, mutations along a DNA
strand: whenever many independent sources each produce a rare event, the events form a \emph{Poisson process}. It is the
continuous-time version of coin tossing, with an exponential waiting time between successes instead of a geometric
one, and it is the basic model of queueing theory and of network traffic.

The exponential law is the only memoryless waiting time, and this one property explains almost everything in the
chapter. We define the Poisson process by its exponential gaps, and prove that it has two other descriptions: by
independent Poisson counts, and by an infinitesimal rate. Given the number of points in an interval, the points are
independent and uniform, which makes thinning and superposition of processes easy. The bus paradox is explained by
size-biased sampling. Le Cam's inequality turns the Poisson limit of Chapter~2 into an explicit error bound, by a
coupling. And the single-server queue M/M/1 settles, when the server is fast enough, into a geometric law whose mean
explodes as the load approaches $100\%$.

\begin{think}
A server receives on average $9$ requests per second and can handle on average $10$. How many requests do you expect to
be waiting or in service at a typical moment: about $1$, about $10$, about $100$? What if it receives $9.9$ per second?
\end{think}

\section{The exponential law}

A random variable $E$ is \emph{exponential with rate $\lambda>0$}, written Exp$(\lambda)$, if $P(E>t)=e^{-\lambda t}$ for $t\ge0$; its density is $\lambda e^{-\lambda t}$ and its mean is $\frac1\lambda$.

\begin{proposition}[memorylessness]\label{prop:memoryless}
Let $T$ be a random variable with $T>0$ and $T<\infty$ almost surely, such that $P(T>s+t)=P(T>s)P(T>t)$ for all $s,t\ge0$; equivalently,
$P(T>s+t\mid T>s)=P(T>t)$ whenever $P(T>s)>0$. Then $T$ is exponential.
\end{proposition}
\begin{proof}
Let $G(t)=P(T>t)$, non-increasing and right-continuous, with $G(0)=1$. If $G(1)=0$ then $G(\frac1n)^n=G(1)=0$, so $G(\frac1n)=0$ for all $n$, contradicting
$G(\frac1n)\to G(0)=1$. If $G(1)=1$ then $G(n)=G(1)^n=1$ for all $n$, so $T=\infty$ almost surely, which is excluded. So $G(1)=e^{-\lambda}$ for some $\lambda>0$.
Then $G(\frac pq)^q=G(p)=G(1)^p$ gives $G(\frac pq)=e^{-\lambda p/q}$ for all positive rationals, and by monotonicity $G(t)=e^{-\lambda t}$ for all $t\ge0$.
\end{proof}

``The bus has not come for ten minutes, so it must be about to arrive'' is false for exponential waiting times: the
remaining wait has the same law as if you had just arrived. The geometric law is the discrete analogue
(Chapter~2).

\begin{lemma}[gaps and arrival times]\label{lem:gamma}
Let $E_1,E_2,\ldots$ be independent Exp$(\lambda)$ and $T_n=E_1+\dots+E_n$. Then $(T_1,\dots,T_n)$ has density $\lambda^ne^{-\lambda t_n}$ on the set
\[
\Delta_n=\{(t_1,\dots,t_n):0<t_1<t_2<\dots<t_n\},
\]
and $0$ elsewhere.
\end{lemma}
\begin{proof}
$(E_1,\dots,E_n)$ has density $\prod_i\lambda e^{-\lambda e_i}=\lambda^ne^{-\lambda(e_1+\dots+e_n)}$ on $]0,\infty[^n$ (independence, Proposition~\ref{prop:joint}). The map $(e_i)\mapsto(t_i)$,
$t_k=e_1+\dots+e_k$, is a linear bijection from $]0,\infty[^n$ onto $\Delta_n$ with triangular matrix of determinant $1$, and $e_1+\dots+e_n=t_n$. Conclude by the
change of variables formula (Proposition~\ref{prop:cov}).
\end{proof}

\section{The Poisson process}

\begin{definition}\label{def:pp}
Let $E_1,E_2,\ldots$ be independent Exp$(\lambda)$, $T_0=0$ and $T_n=E_1+\dots+E_n$. The \emph{Poisson process of rate $\lambda$} is the counting process
\[
N(t)=\#\{n\ge1:T_n\le t\}\qquad(t\ge0).
\]
The $T_n$ are the \emph{arrival times} (or points), and the $E_n$ the \emph{gaps}.
\end{definition}

Since $T_n/n\to\frac1\lambda$ almost surely (Theorem~\ref{thm:slln4}; the gaps have all their moments), $T_n\to\infty$ and $N(t)$ is finite. The path
$t\mapsto N(t)$ is a non-decreasing, right-continuous staircase with jumps of size $1$, and $N(t)\ge n$ if and only if $T_n\le t$.

\begin{theorem}[uniform points]\label{thm:uniform}
Let $t>0$ and $n\ge0$. Then $P(N(t)=n)=e^{-\lambda t}\frac{(\lambda t)^n}{n!}$, and conditionally on $N(t)=n$, the arrival times $(T_1,\dots,T_n)$ have the law of the
order statistics of $n$ independent uniform variables on $[0,t]$: their density is $\frac{n!}{t^n}$ on $\Delta_n(t)=\{0<t_1<\dots<t_n<t\}$.
\end{theorem}
\begin{proof}
Let $B\subseteq\R^n$. The event $\{N(t)=n,\ (T_1,\dots,T_n)\in B\}$ is $\{(T_1,\dots,T_n)\in B,\ T_n\le t<T_n+E_{n+1}\}$. Since $E_{n+1}$ is independent of $(T_1,\dots,T_n)$, with
$P(E_{n+1}>t-t_n)=e^{-\lambda(t-t_n)}$, Lemma~\ref{lem:gamma} and Tonelli show that its probability $q$ is
\[
\begin{aligned}q&=\int_{B\cap\Delta_n(t)}\lambda^ne^{-\lambda t_n}e^{-\lambda(t-t_n)}\dd t_1\cdots\dd t_n\eqbr=\lambda^ne^{-\lambda t}\,|B\cap\Delta_n(t)|,\end{aligned}
\]
where $|\cdot|$ is the volume. With $B=\R^n$ and $|\Delta_n(t)|=\frac{t^n}{n!}$ (the $n!$ orderings of $n$ points of $[0,t]^n$ give congruent pieces, up to the null set where
two coordinates agree), $P(N(t)=n)=e^{-\lambda t}\frac{(\lambda t)^n}{n!}$, and the conditional law has density $\frac{n!}{t^n}$ on $\Delta_n(t)$.

For the order statistics $(U_{(1)},\dots,U_{(n)})$ of independent uniforms $U_1,\dots,U_n$ on $[0,t]$: for $B\subseteq\Delta_n(t)$, $P((U_{(i)})\in B)$ is the sum over the
$n!$ orderings $\sigma$ of $P(U_{\sigma(1)}<\dots<U_{\sigma(n)},\ (U_{\sigma(i)})\in B)=\frac{|B|}{t^n}$, so it is $\frac{n!}{t^n}|B|$: the same law.
\end{proof}

So, given that $n$ points fell in $[0,t]$, they are scattered like $n$ independent uniform points. The counts in disjoint
intervals then follow from the next computation.

\begin{lemma}[colouring]\label{lem:colouring}
Let $N$ be Poisson$(\mu)$, and conditionally on $N=n$ let $(N_1,\dots,N_k)$ be multinomial with $n$ trials and probabilities $p_1,\dots,p_k$ (summing to
$1$). Then $N_1,\dots,N_k$ are independent, and $N_j$ is Poisson$(\mu p_j)$.
\end{lemma}
\begin{proof}
For $n_1+\dots+n_k=n$, the probability $p=P(N_1=n_1,\dots,N_k=n_k)$ is
\[
\begin{aligned}p&=e^{-\mu}\frac{\mu^n}{n!}\cdot\frac{n!}{n_1!\cdots n_k!}p_1^{n_1}\cdots p_k^{n_k}\eqbr=\prod_{j=1}^ke^{-\mu p_j}\frac{(\mu p_j)^{n_j}}{n_j!},\end{aligned}
\]
using $e^{-\mu}=\prod_je^{-\mu p_j}$. A product of the marginal Poisson probabilities is independence.
\end{proof}

\begin{theorem}[independent Poisson increments]\label{thm:increments}
For $0=t_0<t_1<\dots<t_k$, the increments $N(t_j)-N(t_{j-1})$, $1\le j\le k$, are independent, and $N(t_j)-N(t_{j-1})$ is Poisson$(\lambda(t_j-t_{j-1}))$.
\end{theorem}
\begin{proof}
Let $t=t_k$. Conditionally on $N(t)=n$, the points are distributed as $n$ independent uniform points of $[0,t]$, sorted; the increments count
how many of them fall in each interval $]t_{j-1},t_j]$, which does not depend on the sorting, so they are multinomial with probabilities
$\frac{t_j-t_{j-1}}t$. Lemma~\ref{lem:colouring} with $\mu=\lambda t$ concludes.
\end{proof}

The Poisson process has three classical descriptions, and they are equivalent. Call a \emph{counting process} a family $(N(t))_{t\ge0}$ of
random variables with values in $\N$, with $N(0)=0$ and non-decreasing, right-continuous paths.

\begin{theorem}[three definitions]\label{thm:threedefs}
For a counting process $N$ and $\lambda>0$, the following are equivalent.
\begin{enumerate}
\item[(A)] $N$ has the law of the process of Definition~\ref{def:pp}: its jump times $T_n=\inf\{t:N(t)\ge n\}$ have independent Exp$(\lambda)$ gaps.
\item[(B)] For all $0=t_0<t_1<\dots<t_k$, the increments $N(t_j)-N(t_{j-1})$ are independent and Poisson$(\lambda(t_j-t_{j-1}))$.
\item[(C)] The increments over disjoint intervals are independent, the law of $N(t+h)-N(t)$ depends only on $h$, and as $h\to0$,
\[
\begin{gathered}P(N(h)=1)=\lambda h+o(h),\qq P(N(h)\ge2)=o(h).\end{gathered}
\]
\end{enumerate}
\end{theorem}
\begin{proof}
(A)$\Rightarrow$(B) is Theorem~\ref{thm:increments}. (B)$\Rightarrow$(C): $P(N(h)=1)=\lambda he^{-\lambda h}=\lambda h+o(h)$ and $P(N(h)\ge2)=1-e^{-\lambda h}-\lambda he^{-\lambda h}=o(h)$.

(B)$\Rightarrow$(A). By right-continuity, $T_n\le t$ if and only if $N(t)\ge n$. So $P(T_1\le s_1,\dots,T_n\le s_n)$ equals
\[
P\bigl(N(s_1)\ge1,\dots,N(s_n)\ge n\bigr),
\]
and the right side depends only on the joint law of $(N(s_1),\dots,N(s_n))$, which (B) determines (sort the $s_i$ and use the independent
Poisson increments). So $(T_1,\dots,T_n)$ has the same distribution function, hence the same law (Theorem~\ref{thm:uniqueness}, the orthants $]-\infty,s_1]\times\dots\times\,]-\infty,s_n]$
forming a $\pi$-system), as for the process of Definition~\ref{def:pp}, which satisfies (B). Its gaps are therefore independent Exp$(\lambda)$.

(C)$\Rightarrow$(B). Let $p_k(t)=P(N(t)=k)$. Splitting $[0,t+h]$ at $t$ and using independence and stationarity of the increments,
\[
\begin{gathered}p_0(t+h)=p_0(t)p_0(h),\qq p_k(t+h)=p_k(t)p_0(h)+p_{k-1}(t)p_1(h)+r_k(t,h),\end{gathered}
\]
with $0\le r_k(t,h)\le P(N(h)\ge2)=o(h)$. From the first equation, $\frac{p_0(t+h)-p_0(t)}h=-p_0(t)\frac{1-p_0(h)}h\to-\lambda p_0(t)$, since $1-p_0(h)=\lambda h+o(h)$; the same holds for
$h<0$ by writing $p_0(t)=p_0(t-|h|)p_0(|h|)$, and $p_0$ is continuous ($|p_k(t)-p_k(s)|\le P(N(t)-N(s)\ge1)=1-p_0(t-s)$ for $s<t$). So $p_0'=-\lambda p_0$, $p_0(0)=1$: $p_0(t)=e^{-\lambda t}$. In
the same way the second equation gives $p_k'=-\lambda p_k+\lambda p_{k-1}$ with $p_k(0)=0$, so $(e^{\lambda t}p_k)'=\lambda e^{\lambda t}p_{k-1}$, and by induction $p_k(t)=e^{-\lambda t}\frac{(\lambda t)^k}{k!}$. By
stationarity every increment $N(t)-N(s)$ is Poisson$(\lambda(t-s))$, and (C) assumes their independence.
\end{proof}

Description (C) is the one that explains why Poisson processes are everywhere: if in each short time interval an event
happens with probability proportional to its length, independently of the past, then the counts are Poisson.

\section{Thinning and superposition}

Mark each point of a Poisson process independently: given marks $\xi_1,\xi_2,\ldots$, independent of the process, independent, equal to $1$ with
probability $p$ and $2$ with probability $1-p$, let $N_c(t)=\#\{n:T_n\le t,\ \xi_n=c\}$ for $c=1,2$.

\begin{theorem}[thinning]\label{thm:thinning}
$N_1$ and $N_2$ are independent Poisson processes, of rates $\lambda p$ and $\lambda(1-p)$: for all times $0=t_0<\dots<t_k$, the $2k$ increments of $N_1$ and $N_2$ are
independent Poisson variables with parameters $\lambda p(t_j-t_{j-1})$ and $\lambda(1-p)(t_j-t_{j-1})$.
\end{theorem}
\begin{proof}
Let $t=t_k$. Conditionally on $N(t)=n$, the points in $[0,t]$ are the sorted values of $n$ independent uniform variables, and their marks are
independent of them and i.i.d. Attaching the marks in sorted order or in the original order gives the same law of the
(position, mark) pairs, because the marks are i.i.d. and independent of the positions. So the $n$ pairs fall independently in the $2k$
cells (interval $j$, mark $c$), with probabilities $\frac{t_j-t_{j-1}}t\,p_c$ ($p_1=p$, $p_2=1-p$): the cell counts are multinomial. Lemma~\ref{lem:colouring} with $\mu=\lambda t$ concludes.
\end{proof}

\begin{corollary}[superposition]\label{cor:superposition}
If $N_1$ and $N_2$ are independent Poisson processes of rates $\lambda_1$ and $\lambda_2$, then $N_1+N_2$ is a Poisson process of rate $\lambda_1+\lambda_2$. Moreover, the
pair $(N_1,N_2)$ has the law obtained by marking each point of a Poisson process of rate $\lambda_1+\lambda_2$ independently, as $1$ with probability
$\frac{\lambda_1}{\lambda_1+\lambda_2}$ and $2$ otherwise.
\end{corollary}
\begin{proof}
The increments of $N_1+N_2$ over disjoint intervals are independent, and are sums of independent Poisson variables, hence Poisson
with the summed parameters (Chapter~3); so $N_1+N_2$ satisfies (B) of Theorem~\ref{thm:threedefs}. For the second statement, Theorem~\ref{thm:thinning} applied to
a process of rate $\lambda_1+\lambda_2$ produces a pair of independent processes of rates $\lambda_1$ and $\lambda_2$; the joint law of the counts of such a pair is
determined (independent Poisson increments), so it is the law of $(N_1,N_2)$.
\end{proof}

\begin{figure}[ht]
\centering
\fig{ch09-thinning}
\caption{A Poisson process of rate $2$ on $[0,10]$ (top), and its thinning with $p=\frac13$: the kept points (middle, orange) and the
discarded ones (bottom, blue) form two independent Poisson processes, of rates $\frac23$ and $\frac43$. Read from bottom to top, the
figure is a superposition.}
\label{fig:thinning}
\end{figure}

Thinning is how Poisson processes are simulated with a varying rate (keep a point at time $t$ with probability $\frac{\lambda(t)}{\lambda_{\max}}$),
and how network traffic is split between servers. It is also the statement promised in Chapter~3 (Exercise~\ref{ex:thinning}): the
hatched and unhatched eggs are independent Poisson numbers.

\section{The inspection paradox}

Fix a time $t>0$ and look at the gap that contains it: let $A_t=t-T_{N(t)}$ be the \emph{age} (time since the last point, or since $0$) and
$R_t=T_{N(t)+1}-t$ the \emph{residual} time until the next point.

\begin{proposition}\label{prop:inspection}
$R_t$ and $A_t$ are independent, $R_t$ is Exp$(\lambda)$, and $A_t$ has the law of $\min(E,t)$ with $E$ Exp$(\lambda)$. The gap $L_t=A_t+R_t$ containing $t$ has
\[
\mathbb E[L_t]=\frac2\lambda-\frac{e^{-\lambda t}}\lambda\xrightarrow[t\to\infty]{}\frac2\lambda,
\]
twice the mean gap.
\end{proposition}
\begin{proof}
For $s\ge0$ and $0\le a<t$, $\{R_t>s,\ A_t>a\}$ is the event that there is no point in $]t-a,t+s]$, of probability $e^{-\lambda(a+s)}$ by (B). With $a=0$, $R_t$ is
Exp$(\lambda)$; with $s=0$, $P(A_t>a)=e^{-\lambda a}$ for $a<t$, and $A_t\le t$ always, which is the law of $\min(E,t)$. The joint probability is the product, so they
are independent (the events $\{R_t>s\}$ and $\{A_t>a\}$ generate their $\sigma$-algebras; Theorem~\ref{thm:uniqueness}). Finally
$\mathbb E[\min(E,t)]=\int_0^te^{-\lambda a}\dd a=\frac{1-e^{-\lambda t}}\lambda$.
\end{proof}

\begin{figure}[ht]
\centering
\fig{ch09-inspection}
\caption{An observer arriving at a fixed time $t$ (black triangle) lands in a gap; long gaps cover more of the time axis, so they
are more likely to be hit. For Poisson arrivals the observed gap is on average twice as long as a typical one.}
\label{fig:inspection}
\end{figure}

There is no contradiction with the mean gap $\frac1\lambda$: the gap containing a fixed time is not a typical gap, because a long gap is
more likely to contain $t$. Choosing a gap in proportion to its length is called \emph{size-biased sampling}, and it appears whenever one
samples ``by inspection'': the bus you wait for, the class size reported by a student (large classes have more
students to report them), the family size reported by a child, the length of the job a server is busy with when you
look. Exercise~\ref{ex:bus} computes the paradox for any law of gaps: the observed gap has mean $\frac{\mathbb E[X^2]}{\mathbb E[X]}\ge\mathbb E[X]$, with equality only for
regular gaps.

\section{Poisson approximation}

Chapter~2 proved that binomial$(n,\frac\lambda n)$ tends to Poisson$(\lambda)$. A coupling gives an explicit error, for sums of independent Bernoulli
variables with different parameters.

\begin{theorem}[Le Cam's inequality]\label{thm:lecam}
Let $X_1,\dots,X_n$ be independent Bernoulli variables with $P(X_i=1)=p_i$, $S=X_1+\dots+X_n$ and $\mu=p_1+\dots+p_n$. Then
\[
\bigl\|\mathcal L(S)-\text{Poisson}(\mu)\bigr\|\le\sum_{i=1}^np_i^2,
\]
where $\mathcal L(S)$ is the law of $S$ and $\|\cdot\|$ the total variation distance of Chapter~8.
\end{theorem}
\begin{proof}
For one $p\in[0,1]$, build $X$ Bernoulli$(p)$ and $Y$ Poisson$(p)$ from the same uniform $U$ on $[0,1]$: $X=\mathbf 1_{U\ge1-p}$, and $Y=0$ if $U<e^{-p}$, $Y=1$ if
$e^{-p}\le U<e^{-p}(1+p)$, and $Y=k$ on an interval of length $e^{-p}\frac{p^k}{k!}$ further to the right for $k\ge2$. Since $1-p\le e^{-p}$, $X$ and $Y$ differ only when
$1-p\le U<e^{-p}$ or $U\ge e^{-p}(1+p)$, so
\[
\begin{aligned}P(X\ne Y)&=\bigl(e^{-p}-1+p\bigr)+\bigl(1-e^{-p}-pe^{-p}\bigr)\eqbr=p(1-e^{-p})\le p^2.\end{aligned}
\]
Do this independently for each $i$, with independent uniforms $U_i$. Then $S'=Y_1+\dots+Y_n$ is Poisson$(\mu)$, as a sum of independent Poisson variables,
and $P(S\ne S')\le\sum_iP(X_i\ne Y_i)\le\sum_ip_i^2$. The coupling inequality (Lemma~\ref{lem:coupling}) concludes.
\end{proof}

For the binomial$(n,\frac\lambda n)$ the bound is $\frac{\lambda^2}n$. In general $\sum_ip_i^2\le\max_ip_i\cdot\mu$: a sum of many independent rare events, none of which
dominates, is Poisson. For example, if $10^6$ users each click on an advertisement independently, with probabilities
between $10^{-6}$ and $10^{-4}$ adding up to $50$, the number of clicks is within $5\times10^{-3}$ of Poisson$(50)$ in total variation, whatever the
individual probabilities. Le Cam's bound is not sharp; Chen and Stein's method gives $\min(1,\frac1\mu)\sum_ip_i^2$, and also handles
dependent events, such as coincident birthdays.

\section{The M/M/1 queue}

Customers arrive at a single server at the times of a Poisson process of rate $\lambda$ and are served one at a time, in order
of arrival. Service is modelled by an independent Poisson process of rate $\mu$ of \emph{completion signals}: at each signal, if
someone is being served, that customer leaves. (By memorylessness, each service then lasts an Exp$(\mu)$ time; this is the
M/M/1 queue, the letters standing for Markovian arrivals, Markovian services and $1$ server.) Let $Q(t)$ be the number of
customers present at time $t$, waiting or in service, and $\rho=\frac\lambda\mu$, the \emph{load}.

By Corollary~\ref{cor:superposition}, the arrivals and the signals together form a Poisson process of rate $\nu=\lambda+\mu$ whose points are,
independently, arrivals with probability $p=\frac\lambda\nu$ and signals with probability $q=\frac\mu\nu$. So $Q(t)=Y_{K(t)}$, where $K(t)$ is the number of
events by time $t$, a Poisson process of rate $\nu$, and $(Y_m)$ is the Markov chain on $\N$, independent of $K$, with
\[
\begin{gathered}P(k,k+1)=p,\qq P(k,k-1)=q\ (k\ge1),\qq P(0,0)=q,\end{gathered}
\]
started at $Y_0=Q(0)$.

\begin{theorem}[M/M/1]\label{thm:mm1}
If $\rho<1$, then for every initial queue length, $Q(t)$ converges in law, as $t\to\infty$, to the geometric law
\[
\pi(k)=(1-\rho)\rho^k\qquad(k\ge0),
\]
whose mean is $\frac\rho{1-\rho}$. If $\rho\ge1$, the chain $(Y_m)$ has no stationary law.
\end{theorem}
\begin{proof}
\emph{Stationary law.} A stationary law of $(Y_m)$ must balance the flows across each cut between $k$ and $k+1$: for $\pi P=\pi$, summing the
equations for the states $0,\dots,k$ gives $\pi(k)p=\pi(k+1)q$ (the mass entering $\{0,\dots,k\}$ equals the mass leaving it). So $\pi(k)=\pi(0)\rho^k$,
which can be normalised if and only if $\rho<1$, giving $\pi(k)=(1-\rho)\rho^k$; conversely this $\pi$ satisfies detailed balance, hence is stationary
(Proposition~\ref{prop:detailed}).

\emph{Convergence of $(Y_m)$, by a monotone coupling.} Drive two copies $Y$ and $Y'$, started at $x$ and from $\pi$, with the same sequence of
independent signs $\epsilon_m$ ($+1$ with probability $p$): $Y_{m+1}=\max(Y_m+\epsilon_{m+1},0)$, and likewise $Y'$. Each copy is a chain with matrix $P$. The map
$y\mapsto\max(y+\epsilon,0)$ is non-decreasing and does not increase distances, so the order between $Y_m$ and $Y'_m$ never changes, and once the upper
one is at $0$ both are at $0$ and stay together. The upper copy is the chain started at $\max(x,Y'_0)$; before it reaches $0$ it is a walk
with $p<q$, which reaches $0$ almost surely (Corollary~\ref{cor:infinitebank}). So $P(Y_m\ne Y'_m)\to0$, and by the coupling inequality
$\|\mathcal L(Y_m)-\pi\|\le P(Y_m\ne Y'_m)\to0$.

\emph{Continuous time.} As $K(t)$ is independent of $Y$, $P(Q(t)=k)=\sum_mP(K(t)=m)P(Y_m=k)$, so $\delta(t)=\|\mathcal L(Q(t))-\pi\|$ satisfies
\[
\begin{aligned}\delta(t)&\le\sum_mP(K(t)=m)\,\|\mathcal L(Y_m)-\pi\|\eqbr\le P(K(t)<M)+\sup_{m\ge M}\|\mathcal L(Y_m)-\pi\|.\end{aligned}
\]
For fixed $M$, $P(K(t)<M)\to0$ as $t\to\infty$; then let $M\to\infty$. The mean of $\pi$ is $\sum_kk(1-\rho)\rho^k=\frac\rho{1-\rho}$.
\end{proof}

\begin{figure}[ht]
\centering
\fig{ch09-mm1}
\caption{The number of customers in an M/M/1 queue with $\mu=1$, for $\lambda=0.5$ (blue) and $\lambda=0.9$ (orange), over $300$ time units. The
stationary means are $1$ and $9$; at high load the queue makes long excursions.}
\label{fig:mm1}
\end{figure}

\begin{beauty}
The mean queue length $\frac\rho{1-\rho}$ is $1$ at load $50\%$, $9$ at $90\%$ and $99$ at $99\%$. A system that is ``only'' $10\%$ below saturation already has
long queues, and the last few percent of capacity are paid for in waiting time: this one formula is behind the rule of
thumb that servers, roads and hospitals should not run near full load.
\end{beauty}

\section{Exercises}

\begin{exercise}[$\star$]\label{ex:race}
\emph{Competing exponentials.} Let $E_1,\dots,E_k$ be independent, $E_i$ Exp$(\lambda_i)$, $\Lambda=\lambda_1+\dots+\lambda_k$, $M=\min_iE_i$ and $I$ the index achieving it. Show that $M$ is
Exp$(\Lambda)$, that $P(I=i)=\frac{\lambda_i}\Lambda$, and that $M$ and $I$ are independent.
\end{exercise}

\begin{exercise}[$\star$]\label{ex:binomialsplit}
For a Poisson process, show that conditionally on $N(t)=n$, $N(s)$ is binomial$(n,\frac st)$ for $0<s<t$. Passengers arrive at a bus stop as a
Poisson process of rate $\lambda$, and a bus leaves at time $t$: show that their expected total waiting time is $\frac{\lambda t^2}2$.
\end{exercise}

\begin{exercise}[$\star\star$]\label{ex:records2}
Let $R$ be the number of records among $n$ numbers in random order that occur at positions $m+1,\dots,n$ (Exercise~\ref{ex:records}).
Show that $R$ is within $\frac1m$ of Poisson$(H_n-H_m)$ in total variation.
\end{exercise}

\begin{exercise}[$\star\star$]\label{ex:coupon2}
\emph{Coupons in continuous time.} Coupons of $n$ types arrive at the times of a Poisson process of rate $1$, each of a uniformly random
type. Let $T$ be the time when every type has arrived. Show that $P(T\le t)=(1-e^{-t/n})^n$, that $\mathbb E[T]=nH_n$, and that
$P(T\le n\ln n+cn)\to e^{-e^{-c}}$ for every real $c$ (the Gumbel law).
\end{exercise}

\begin{exercise}[$\star\star$]\label{ex:little}
In a stationary M/M/1 queue, a customer who arrives finds a number of customers distributed according to $\pi$ (take this
from the event chain $(Y_m)$: an arrival is an event chosen independently of the current state). Show that the customer's
expected time in the system is $\frac1{\mu-\lambda}$, and check \emph{Little's law} $L=\lambda W$, where $L$ is the mean number in the system and $W$ the mean
time spent in it.
\end{exercise}

\begin{exercise}[$\star\star$]\label{ex:bus}
\emph{The bus paradox for any gaps.} Buses arrive at times $T_n=X_1+\dots+X_n$, with independent gaps of the same law, bounded by a
constant and with $\mathbb E[X]>0$. Let $L_t$ be the length of the gap containing $t$. Show that $\frac1{T_n}\int_0^{T_n}L_t\dd t\to\frac{\mathbb E[X^2]}{\mathbb E[X]}$ almost surely. Compute this for
gaps uniform on $[0,2m]$ and for constant gaps $m$.
\end{exercise}

\begin{exercise}[$\star$]\label{ex:superpose}
Two independent Poisson processes of rates $\lambda_1$ and $\lambda_2$ are observed. Show that the number of points of the first one before the first
point of the second is geometric: $P(k\text{ points})=\bigl(\frac{\lambda_1}{\lambda_1+\lambda_2}\bigr)^k\frac{\lambda_2}{\lambda_1+\lambda_2}$.
\end{exercise}

\section{Hints}

\begin{hint}{ex:race}
$P(M>t,\ I=i)=\int_t^\infty\lambda_ie^{-\lambda_is}\prod_{j\ne i}e^{-\lambda_js}\dd s$.
\end{hint}
\begin{hint}{ex:binomialsplit}
Theorem~\ref{thm:uniform}. Given $N(t)=n$, the total waiting time is $\sum_i(t-U_i)$ with $U_i$ uniform.
\end{hint}
\begin{hint}{ex:records2}
The record indicators are independent Bernoulli$(\frac1k)$, and $\sum_{k>m}\frac1{k^2}\le\frac1m$.
\end{hint}
\begin{hint}{ex:coupon2}
Thin the arrival process by types: $n$ independent Poisson processes of rate $\frac1n$. For the mean, write $T=\sum_{j=1}^K E_j$ with $K$ the discrete
coupon collector's time, independent of the gaps $E_j$.
\end{hint}
\begin{hint}{ex:little}
A customer who finds $k$ customers waits for $k+1$ services; the one in progress has a residual Exp$(\mu)$ duration by memorylessness.
\end{hint}
\begin{hint}{ex:bus}
$\int_0^{T_n}L_t\dd t=\sum_{i\le n}X_i^2$. Use the strong law for the $X_i$ and for the $X_i^2$.
\end{hint}
\begin{hint}{ex:superpose}
Corollary~\ref{cor:superposition}: the points of the superposition are marked independently.
\end{hint}

\section{Solutions}

\begin{solution}{ex:race}
\idea Each clock rings at its own rate; the first ring comes at the total rate, and which clock rings is decided by the
rates alone.
\writeup By independence,
\[
\begin{aligned}P(M>t,\ I=i)&=\int_t^\infty\lambda_ie^{-\lambda_is}\prod_{j\ne i}e^{-\lambda_js}\dd s\eqbr=\frac{\lambda_i}\Lambda e^{-\Lambda t},\end{aligned}
\]
since $\{M>t,\ I=i\}=\{E_i>t,\ E_j>E_i\ \forall j\ne i\}$ (ties have probability $0$) and $P(E_j>s)=e^{-\lambda_js}$ (Tonelli). With $t=0$, $P(I=i)=\frac{\lambda_i}\Lambda$; summing over $i$,
$P(M>t)=e^{-\Lambda t}$. The joint probability is the product, so $M$ and $I$ are independent.
\end{solution}

\begin{solution}{ex:binomialsplit}
\idea Given the number of points, they are independent uniform points, and each falls before $s$ with probability $\frac st$.
\writeup By Theorem~\ref{thm:uniform}, given $N(t)=n$ the points have the law of $n$ sorted independent uniforms on $[0,t]$, and $N(s)$ counts those below $s$:
binomial$(n,\frac st)$. The passengers arriving in $[0,t]$ wait in total $\sum_{i\le N(t)}(t-T_i)$. Given $N(t)=n$, this is $\sum_{i\le n}(t-U_i)$ in law (the sum does not
depend on the order), of mean $\frac{nt}2$. So the expected total is $\mathbb E[N(t)]\frac t2=\frac{\lambda t^2}2$.
\end{solution}

\begin{solution}{ex:records2}
\idea Independent rare events with known probabilities: exactly Le Cam's setting.
\writeup By Exercise~\ref{ex:records}, the indicators $I_k$ that position $k$ holds a record are independent Bernoulli$(\frac1k)$. So $R=\sum_{k=m+1}^nI_k$ and, by
Theorem~\ref{thm:lecam}, its distance to Poisson$(\sum_{k>m}\frac1k)=$ Poisson$(H_n-H_m)$ is at most $\sum_{k>m}\frac1{k^2}\le\sum_{k>m}\frac1{k(k-1)}=\frac1m$ (a telescoping sum). (For all
the records, $m=1$, the bound $1$ is useless, because the first positions are records with probabilities $1$, $\frac12$, \dots, not rare.)
\end{solution}

\begin{solution}{ex:coupon2}
\idea Poissonisation: in continuous time the types arrive independently, so ``all types present'' is a product of
independent events.
\writeup Marking each arrival by its type (probability $\frac1n$ each) and applying Theorem~\ref{thm:thinning} repeatedly (or Lemma~\ref{lem:colouring} with $n$ colours),
the arrivals of the $n$ types are independent Poisson processes of rate $\frac1n$. Type $j$ first arrives at an Exp$(\frac1n)$ time $F_j$, independent
across types, and $T=\max_jF_j$. So $P(T\le t)=\prod_jP(F_j\le t)=(1-e^{-t/n})^n$. The type sequence is independent of the gaps $E_1,E_2,\ldots$, so $T=E_1+\dots+E_K$, where
$K$ is the discrete coupon collector's time (Theorem~\ref{thm:coupon}), independent of the $E_j$; by Wald's identity (Corollary~\ref{cor:wald}),
$\mathbb E[T]=\mathbb E[K]\,\mathbb E[E_1]=nH_n$. Finally, with $t=n\ln n+cn$, $e^{-t/n}=\frac{e^{-c}}n$ and $(1-\frac{e^{-c}}n)^n\to e^{-e^{-c}}$.
\end{solution}

\begin{solution}{ex:little}
\idea The newcomer waits for everyone ahead, and memorylessness makes each of those services, including the one in
progress, a fresh Exp$(\mu)$ time.
\writeup Given that the customer finds $k$ customers, the time in the system is the residual service of the customer in service (Exp$(\mu)$ by
memorylessness, if $k\ge1$), plus $k-1$ full services, plus the customer's own: $k+1$ independent Exp$(\mu)$ times, of mean $\frac{k+1}\mu$. Averaging over
$k\sim\pi$, $W=\frac{\mathbb E[k]+1}\mu=\frac1\mu\bigl(\frac\rho{1-\rho}+1\bigr)=\frac1{\mu(1-\rho)}=\frac1{\mu-\lambda}$. Then $\lambda W=\frac\lambda{\mu-\lambda}=\frac\rho{1-\rho}=L$. Little's law holds for very general queues,
with essentially the same reason: count customer-seconds in two ways.
\end{solution}

\begin{solution}{ex:bus}
\idea Integrate over time: each gap of length $X_i$ contributes $X_i$ times its own length.
\writeup On the $i$-th gap, of length $X_i$, $L_t=X_i$, so $\int_0^{T_n}L_t\dd t=\sum_{i\le n}X_i^2$ (the gaps of length $0$ contribute nothing). By the strong law (Theorem~\ref{thm:slln4},
all moments being finite for bounded gaps), $\frac1n\sum X_i^2\to\mathbb E[X^2]$ and $\frac{T_n}n\to\mathbb E[X]>0$ almost surely, so the ratio tends to $\frac{\mathbb E[X^2]}{\mathbb E[X]}$. This is at
least $\mathbb E[X]$ by Cauchy--Schwarz ($\mathbb E[X]^2\le\mathbb E[X^2]$), with equality only for constant gaps. For gaps uniform on $[0,2m]$: $\mathbb E[X^2]=\frac{4m^2}3$, giving $\frac{4m}3$; for
constant gaps, $m$. For exponential gaps of mean $m$, $\mathbb E[X^2]=2m^2$ and the answer is $2m$, as in Proposition~\ref{prop:inspection}.
\end{solution}

\begin{solution}{ex:superpose}
\idea In the superposition, each point is a coin flip between the two processes, so we wait for the first ``tail''.
\writeup By Corollary~\ref{cor:superposition}, the points of $N_1+N_2$ are marked independently, $1$ with probability $r=\frac{\lambda_1}{\lambda_1+\lambda_2}$. The number of points of $N_1$ before
the first point of $N_2$ is the number of marks $1$ before the first mark $2$, which is $k$ with probability $r^k(1-r)$.
\end{solution}

\section*{Sources and further reading}

\begin{itemize}
\item G.~Grimmett and D.~Stirzaker, \emph{Probability and Random Processes}, 4th ed., Oxford University Press, 2020,
Chapter~6: the Poisson process, its three definitions and the inspection paradox.
\item J.~F.~C.~Kingman, \emph{Poisson Processes}, Oxford University Press, 1993. Colouring, thinning and superposition,
in a short and elegant book.
\item L.~Le Cam, \emph{An approximation theorem for the Poisson binomial distribution}, Pacific Journal of Mathematics
10 (1960), 1181--1197.
\item A.~D.~Barbour, L.~Holst and S.~Janson, \emph{Poisson Approximation}, Oxford University Press, 1992. The Chen--Stein
method.
\item L.~Kleinrock, \emph{Queueing Systems}, Vol.~1, Wiley, 1975. The M/M/1 queue and its relatives.
\item J.~D.~C.~Little, \emph{A proof for the queuing formula $L=\lambda W$}, Operations Research 9 (1961), 383--387.
\item S.~L.~Feld, \emph{Why your friends have more friends than you do}, American Journal of Sociology 96 (1991),
1464--1477. Size-biased sampling in social networks.
\item L.~von Bortkiewicz, \emph{Das Gesetz der kleinen Zahlen}, Teubner, Leipzig, 1898. The horse-kick table.
\end{itemize}
